% Part: proof-theory
% Chapter: natural-deduction
% Section: introduction

\documentclass[../../../include/open-logic-section]{subfiles}

\begin{document}

\olfileid{pt}{nat}{int}

\olsection{పరిచయం}

సహజ నిగమనం అనేది \tetoken{నిరూపణ}{!!{proof}} పద్ధతుల కుటుంబం. ఇందులో ప్రతి తార్కిక సంయోజకానికి లేదా పరిమాణీకరణికానికి ఒక జత నిగమనాలు ఉంటాయి: ఆ సంయోజకాన్ని ``ప్రవేశపెట్టే'' నిగమనం (\tetoken{ఒక ప్రవేశ}{!!a{introduction}} నియమం), దాన్ని ``తొలగించే'' నిగమనం (\tetoken{ఒక తొలగింపు}{!!a{elimination}} నియమం). ఉదాహరణకు, \Intro{\lif} నిగమనాల ఫలితాలు~$!A \lif !B$ రూపపు \tetoken{సూత్రాలు}{!!{formula}s}; \Elim{\lif} నిగమనాల పూర్వపక్షాలలో~$!A \lif !B$ రూపపు \tetoken{సూత్రాలు}{!!{formula}s} ఉంటాయి.

సహజ నిగమనం యొక్క మొదటి రెండు రూపాలను 1930లలో గెర్హార్డ్ గెంట్జెన్, స్టానిస్వావ్ యాష్కోవ్‌స్కీ ప్రవేశపెట్టారు. నిరూపణ సిద్ధాంతవేత్తలు ఎక్కువగా గెంట్జెన్ పద్ధతినే చర్చిస్తారు; బోధనలో విస్తృతంగా వాడే పద్ధతులకు యాష్కోవ్‌స్కీ పద్ధతి మూలం. రెండు పద్ధతుల్లోనూ \emph{ఉపపత్తులు} నుంచి ప్రారంభించవచ్చు. ఇవి \tetoken{నిరూపించ}{!!{prove}}నవసరం లేని \tetoken{సూత్రాలు}{!!{formula}s}, కానీ ఇతర \tetoken{సూత్రాలను}{!!{formula}s} \tetoken{నిరూపించడానికి}{!!{proving}} ఉపయోగపడతాయి. మొత్తం \tetoken{నిరూపణ}{!!{proof}} అలాంటి ఉపపత్తులపై ఆధారపడవచ్చు. అవి నిరూపించబడలేదు కనుక ఆ \tetoken{నిరూపణ}{!!{proof}} తన ఫలితాన్ని (అంటే దాని \emph{చివరి-\tetoken{సూత్రాన్ని}{!!{formula}}}) నేరుగా స్థాపించదు; ఉపపత్తుల నుంచి ఆ ఫలితం అనుగమిస్తుందని మాత్రమే స్థాపిస్తుంది.

కొన్ని నిగమన నియమాలలో, వాటి ఫలితాలు కొన్ని ఉపపత్తులపై ఇక ఆధారపడవు: అవి ఆ ఉపపత్తులను ``విసర్జిస్తాయి.'' ఉదాహరణకు, \Intro{\lif} నియమం~$!A$ ఉపపత్తి నుంచి~$!B$ ఫలితానికి గల \tetoken{ఒక నిరూపణను}{!!a{proof}} తీసుకుని, $!A \lif !B$ను ఫలితంగా ఇస్తుంది. $!A \lif !B$తో ముగిసే \tetoken{నిరూపణ}{!!{proof}} ఇక~$!A$పై ఆధారపడదు; $!A$ ఉపపత్తి విసర్జించబడింది. గెంట్జెన్ పద్ధతులలో \tetoken{నిరూపణలు}{!!{proof}s} \tetoken{సూత్రాల}{!!{formula}s} వృక్షాలు; ఉపపత్తులు అత్యుపరి \tetoken{సూత్రాలు}{!!{formula}s};~\Intro{\lif} వంటి నియమాలపై ఏ ఉపపత్తులు విసర్జించబడ్డాయో తెలిపే గుర్తులు ఉంటాయి. యాష్కోవ్‌స్కీ పద్ధతులలో, ఒక ఉపపత్తిపై ఆధారపడే \tetoken{నిరూపణ}{!!{proof}} భాగాన్ని ఏదో విధంగా వేరుచేస్తారు; ఉదాహరణకు, నిలువు గీతకు అవతల లోపలికి జరిపి లేదా ఒక పెట్టెలో ఉంచుతారు. ఆ పెట్టె మొదటి వరుసలో ఉపపత్తి~$!A$, చివరి వరుసలో షరతు వాక్యంలోని ఫలభాగం~$!B$ ఉంటాయి; పెట్టె బయట \Intro{\lif} ఫలితం~$!A \lif !B$ ఉంటుంది, ఆ ఉపపత్తి విసర్జించబడుతుంది.
\sourcecorrection{OLTEPTNATINT-001}{యాష్కోవ్‌స్కీ పెట్టెలో ఉపపత్తి $!A$ను విసర్జిస్తామని చెబుతూ, మూలం ఫలభాగం $!B$నే మొదటి, చివరి వరుసల రెండింటిలో ఉంచింది. పెట్టె మొదటిలో $!A$, చివరలో $!B$ ఉండేలా సరిచేశాం.}

ఈ పద్ధతులను ``సహజ''ం అంటారు, ఎందుకంటే వాటి నిగమన నియమాలు మనం (కనీసం గణితవేత్తలు) సహజంగా తర్కించే తీరును ప్రతిబింబిస్తాయి. ఒక షరతుతో కూడిన ఫలితాన్ని స్థాపించడానికి, ముందుభాగాన్ని ఉపపత్తిగా తీసుకుని దాని నుంచి ఫలభాగాన్ని నిరూపిస్తాం. అదే~\Intro{\lif}. షరతు వాక్యం $!A \lif !B$, దాని ముందుభాగం~$!A$ తెలిసినప్పుడు~$!B$ను నిర్ధారిస్తాం. అదే~\Elim{\lif}. వియోజనం~$!A \lor !B$ను ఉపపత్తిగా తీసుకున్నప్పుడు ఒక ఫలితం వస్తుందని చూపడానికి సందర్భాలవారీ నిరూపణ వాడతాం. అదే~\Elim{\lor}. సాధారణ వాదన~$\lforall[x][!A(x)]$ను నిరూపించడానికి, ఏదైనా స్వేచ్ఛగా ఎంచుకున్న~$c$ కోసం $!A(c)$ నిజమని నిరూపిస్తాం. అదే~\Intro{\lforall}. ఇలాగే మిగిలినవి.

ఉదాహరణకు గెంట్జెన్-శైలి సహజ నిగమనంలో~$!A \lif (!A \lor !B)$కు ఒక సరళమైన \tetoken{నిరూపణ}{!!{proof}} ఇది:
\[
\AxiomC{$\Discharge{!A}{x}$}
\RightLabel{\Intro{\lor}}
\UnaryInfC{$!A \lor !B$}
\DischargeRule{\Intro{\lif}}{x}
\UnaryInfC{$!A \lif (!A \lor !B)$}
\DisplayProof
\]
ఇది~$!A$ ఉపపత్తితో మొదలై, \Intro{\lor}తో~$!A
\lor !B$ను, తరువాత \Intro{\lif}తో~$!A \lif (!A \lor !B)$ను నిగమిస్తుంది. \Intro{\lif} నిగమనం~$!A$ ఉపపత్తిని విసర్జిస్తుంది; దానినే~$x$ గుర్తు సూచిస్తుంది.

నిరూపణ సిద్ధాంతవేత్తలు \tetoken{సూత్రాల}{!!{formula}s} వృక్షాలపై పనిచేసే గెంట్జెన్ అసలు పద్ధతులను ఇష్టపడతారు; అయితే వాటి ఒక రూపాంతరం కూడా నిరూపణ సిద్ధాంతంలో ముఖ్యమైనది. ఉపపత్తులు~$\Gamma$ నుంచి~$!A$కు గల \tetoken{ఒక నిరూపణ}{!!^a{proof}} ఉంటే,~$\Gamma$ నుంచి~$!A$ అనుగమిస్తుందని, అంటే $\Gamma \Entails !A$ అని తెలుస్తుంది. సహజ నిగమన నియమాలను అలాంటి అనుగమనాల గురించి చెప్పే నియమాలుగా చదవవచ్చు. ఉదాహరణకు, \Intro{\lif} ప్రకారం $\Gamma + !A
\Entails !B$ అయితే $\Gamma \Entails !A \lif !B$. అందువల్ల నిగమన నియమాల పూర్వపక్షాలూ ఫలితమూ ఉపపత్తులు, వాటి నుంచి వచ్చే ఫలితం రెండింటిపైనా పనిచేసేలా వాటిని నేరుగా రూపొందించడం సహజం; అంటే అవి $\Gamma \Sequent !A$ సీక్వెంట్‌లపై పనిచేస్తాయి. పై సరళ నిరూపణ అటువంటి పద్ధతిలో ఇలా ఉంటుంది:
\[
\Axiom$x:!A \fCenter !A$
\RightLabel{\Intro{\lor}}
\UnaryInf$x:!A \fCenter !A \lor !B$
\DischargeRule{\Intro{\lif}}{x}
\UnaryInf$\fCenter !A \lif (!A \lor !B)$
\DisplayProof
\]
చూసినట్లుగా నియమాలు అవే: అవి~$\Sequent$కు కుడివైపు ఉన్న \tetoken{సూత్రంపై}{!!{formula}} పనిచేస్తాయి; ఎడమవైపు ఉన్న \tetoken{సూత్రాలు}{!!{formula}s} ప్రతి దశలో అది ఆధారపడే ఉపపత్తులను మాత్రమే జాబితా చేస్తాయి.
\sourcecorrection{OLTEPTNATINT-002}{మూలం \Intro{\lif} ద్వారా $\Gamma + !A \Entails !B$ నుంచి $\Gamma \Entails !B$ వస్తుందని చెప్పింది. ఇక్కడ సరైన ఫలితం $\Gamma \Entails !A \lif !B$; దాన్ని సరిచేశాం.}

\tetoken{ప్రవేశ}{!!{introduction}}/\tetoken{తొలగింపు}{!!{elimination}} నియమ జతల సంపూర్ణ సమితి ఒక్కటే సాంప్రదాయిక తర్కానికి సంపూర్ణం కాదు. $\lfalse$కు సంబంధించిన మరో రెండు నియమాలు కావాలి. మొదటిది ``అంతఃప్రజ్ఞావాద అసంభవ నియమం''~\FalseInt, లేదా ``విస్ఫోటనం''; $\lfalse$ నుంచి దేనినైనా నిగమించవచ్చని ఇది చెబుతుంది. రెండోది పరోక్ష నిరూపణకు సంబంధించిన సాంప్రదాయిక సూత్రం~\FalseCl{} (లేదా ``సాంప్రదాయిక అసంభవ నియమం'');~$\lnot !A$ నుంచి~$\lfalse$ను \tetoken{నిరూపించగలిగితే}{!!{prove}},~$!A$ \tetoken{నిరూపించబడిందని}{!!{prove}} అది చెబుతుంది. \FalseClను తీసివేస్తే అంతఃప్రజ్ఞావాద తర్కానికి సరిపోయే పద్ధతి వస్తుంది. రెండింటినీ తీసివేస్తే దాన్ని ``కనిష్ఠ తర్కం'' అంటారు.

సాంప్రదాయిక, అంతఃప్రజ్ఞావాద తర్కాల కోసం గెంట్జెన్-శైలి సహజ నిగమనాన్ని చర్చిస్తాం. అవి \Log{N1c}, \Log{N1i} పద్ధతులు. వాటికి అనురూపంగా, ఒక్కొక్క \tetoken{సూత్రం}{!!{formula}s}~$!A$కు బదులుగా $\Gamma \Sequent !A$ సీక్వెంట్‌లపై పనిచేసేవి \Log{N2c}, \Log{N2i} పద్ధతులు.

\end{document}
